% Shared visual system for every TikZ figure in this report.
% Colours are sampled from repo/layout.png, the annotated layout image that
% ships with the puzzle; see figures/src/palette.py for the measurement.
\usetikzlibrary{arrows.meta,positioning,calc,fit,backgrounds,shapes.geometric,
  decorations.pathreplacing,patterns}

\definecolor{figink}{HTML}{172033}     % body ink
\definecolor{figsoft}{HTML}{5B6472}    % secondary text
\definecolor{figaccent}{HTML}{3157A4}  % structure, arrows, frames
\definecolor{figblue}{HTML}{2196F3}    % layout.png: routing
\definecolor{figviolet}{HTML}{B388FF}  % layout.png: straps
\definecolor{figorange}{HTML}{FF8C1A}  % layout.png: vertical bars
\definecolor{figmagenta}{HTML}{E040E0} % layout.png: fine routes
\definecolor{figgold}{HTML}{C9A227}    % layout.png: local interconnect
\definecolor{figgreen}{HTML}{0A7D33}   % layout.png: output labels; used for
                                       % verified facts and final results
% The layout violet is a fill colour: as a hairline stroke or as small text it
% reads as grey. This is the same hue carried toward the ink so that clock
% wires and clock labels stay legible at figure sizes.
\definecolor{figclkline}{HTML}{7B52C4}
\definecolor{figrule}{HTML}{B9C0CC}    % muted outlines
\definecolor{figmute}{HTML}{F1F4F8}    % muted fills
\definecolor{figwash}{HTML}{EEF3FC}    % pale accent wash

% Narrow boxes hyphenate badly at this size, and a broken word costs more
% legibility than a ragged edge saves.
\newcommand{\fignohyphen}{\hyphenchar\font=-1 \hyphenpenalty=10000
  \exhyphenpenalty=10000 \tolerance=3000 \emergencystretch=3em}

% A card: eyebrow, title, body, hairline, and a monospace footer pinned to the
% bottom edge. Keeping the footer inside the frame stops the artefact names
% floating loose underneath, and a fixed height keeps a row of cards even.
% #1 name  #2 (x,y) of the top-left corner  #3 width/cm  #4 height/cm
% #5 style  #6 eyebrow  #7 title  #8 body  #9 footer
% The footer is anchored to the bottom edge and grows upward, so a two- or
% three-line footer stays inside the frame instead of spilling past it, and the
% hairline is placed from the footer's own top edge rather than a guessed offset.
\newcommand{\figcard}[9]{%
  \node[#5, anchor=north west, minimum width=#3cm, minimum height=#4cm,
        inner sep=0pt] (#1) at #2 {};
  \node[ftag, anchor=north west] at ($(#1.north west)+(0.22,-0.13)$) {#6};
  \node[figfont, anchor=north west, font=\sffamily\scriptsize\bfseries]
    at ($(#1.north west)+(0.22,-0.42)$) {#7};
  \node[figfont, anchor=north west, align=left,
        text width={\dimexpr #3cm-0.44cm}]
    at ($(#1.north west)+(0.22,-0.76)$) {\fignohyphen #8};
  \node[fmonosoft, anchor=south west, align=left,
        text width={\dimexpr #3cm-0.44cm}]
    (#1-foot) at ($(#1.south west)+(0.22,0.14)$) {#9};
  \draw[figrule, line width=0.4pt]
    ($(#1-foot.north west)+(0,0.10)$) -- ($(#1-foot.north east)+(0,0.10)$);
}

\tikzset{
  % applies inside every node's own text group, where it still governs how
  % lines are broken
  every node/.append style={execute at begin node=\fignohyphen},
  figfont/.style={font=\sffamily\scriptsize, text=figink},
  % a plain content box
  fbox/.style={figfont, draw=figaccent, fill=white, rounded corners=2pt,
    line width=0.5pt, align=center, inner sep=4pt},
  % a stage in a pipeline: title line plus supporting lines
  fstage/.style={fbox, fill=figwash, align=left, text width=#1},
  fstage/.default=3.1cm,
  % something established by evidence
  fgood/.style={fbox, draw=figgreen, fill=figgreen!5},
  % something taken on trust from outside the file
  ftrust/.style={fbox, draw=figorange, fill=figorange!6},
  % de-emphasised context
  fmuted/.style={figfont, draw=figrule, fill=figmute, rounded corners=2pt,
    line width=0.4pt, align=center, inner sep=3.5pt, text=figsoft},
  % arrows
  fflow/.style={-{Latex[length=2.1mm,width=1.5mm]}, draw=figaccent,
    line width=0.6pt},
  fthin/.style={-{Latex[length=1.6mm,width=1.2mm]}, draw=figsoft,
    line width=0.4pt},
  fdep/.style={draw=figsoft, line width=0.4pt, dash pattern=on 1.6pt off 1.4pt},
  % labels
  ftag/.style={font=\sffamily\tiny, text=figsoft, inner sep=1.5pt},
  fmono/.style={font=\ttfamily\tiny, text=figink, inner sep=1pt},
  fmonosoft/.style={font=\ttfamily\tiny, text=figsoft, inner sep=1pt},
  fhead/.style={font=\sffamily\scriptsize\bfseries, text=figaccent},
  %
  % ---- signal roles -----------------------------------------------------
  % One colour per role, used the same way in every schematic figure, so a
  % reader who learns the key once can carry it between diagrams:
  %   data      routing blue    the value being computed or carried
  %   clock     violet          when state advances, never a data term
  %   enable    magenta         whether this edge captures anything
  %   reset     orange          the asynchronous override
  % Green stays reserved for checked facts and final results. Captions identify
  % the smaller subset checked against genuinely independent evidence.
  fdata/.style={-{Latex[length=2.1mm,width=1.5mm]}, draw=figblue,
    line width=0.9pt},
  fclk/.style={-{Latex[length=2.1mm,width=1.5mm]}, draw=figclkline,
    line width=0.9pt},
  fen/.style={-{Latex[length=2.1mm,width=1.5mm]}, draw=figmagenta,
    line width=0.9pt},
  frst/.style={-{Latex[length=2.1mm,width=1.5mm]}, draw=figorange,
    line width=0.9pt},
  fverified/.style={-{Latex[length=2.1mm,width=1.5mm]}, draw=figgreen,
    line width=0.9pt},
  % a bus, drawn heavier than a single wire
  fbus/.style={-{Latex[length=2.4mm,width=2.0mm]}, draw=figblue,
    line width=1.7pt},
  % an explanatory strip that runs under a figure
  fnote/.style={fmuted, align=left, inner sep=6pt},
  % a small square swatch used in figure keys
  fswatch/.style={draw=#1, fill=#1, line width=0pt},
}

% A key entry: a short coloured rule plus a label, for in-figure legends.
% #1 colour  #2 label
\newcommand{\figkeyitem}[2]{%
  \tikz[baseline=-0.5ex]{\draw[#1, line width=0.9pt] (0,0) -- (0.34,0);}\,%
  {\ttfamily\tiny #2}}

% A section heading inside a figure, with a hairline under it.
% #1 node name  #2 (x,y)  #3 width/cm  #4 text
\newcommand{\figband}[4]{%
  \node[fhead, anchor=north west, inner sep=1pt] (#1) at #2 {#4};
  \draw[figrule, line width=0.4pt]
    ($(#1.south west)+(0,-0.02)$) -- ++(#3,0);
}
